\subsection{Nilradical and minimal prime ideals}

In a (commutative) ring A the set of nilpotent elements is an ideal, for if x, y are elements of A such that \(x^{m} = y^{n} = 0\) , then \((x + y)^{m+n} = 0\) by the binomial theorem.

DEFINITION 4. The ideal of nilpotent elements of a (commutative) ring A is called the nilradical of A. If a is an ideal of A, the inverse image, under the canonical mapping \(A \rightarrow A/a\) , of the nilradical of A/a is called the radical of a.

We often denote by \(\mathbf{r}(\mathfrak{a})\) the radical of an ideal a of A.

To say that an element \(x \in A\) belongs to the radical of \(a\) means therefore that there exists an integer \(n > 0\) such that \(x^n \in a\) . If \(f\) is a homomorphism from \(A\) to a ring \(B\) and \(b\) is an ideal of \(B\) , the radical of \(f(b)\) is the inverse image under \(f\) of the radical of \(b\) , for to say that \(x^n \in f(b)\) means that \((f(x))^n \in b\) .

The nilradical of a ring A is contained in its Jacobson radical (Algebra, Chapter VIII, § 6, no. 3, Corollary 3 to Theorem 1) but may be distinct from it; it is always equal to it if A is Artinian (Algebra, Chapter VIII, § 6, no. 4, Theorem 3).

We say that a prime ideal p of a ring A is a minimal prime ideal if it is minimal in the set of prime ideals of A ordered by inclusion.

PROPOSITION 12. Let \(\mathfrak{p}\) be a minimal prime ideal of a ring \(\mathbf{A}\) . For all \(x \in \mathfrak{p}\) , there exists \(s \in \mathbf{A} - \mathfrak{p}\) and an integer \(k > 0\) such that \(sx^k = 0\) .

The set S of elements of the form \(sx^{k}\) (k an integer \textgreater0, \(s \in A - p\) ) is a multiplicative subset of A. If \(0 \notin S\) , there would exist a prime ideal \(p'\) not meeting S (no. 5, Corollary 2 to Proposition 11). Then \(p' \subset p\) and \(p' \neq p\) since \(x \notin p'\) , contrary to the hypothesis that p is minimal.

PROPOSITION 13. The nilradical of a ring A is the intersection of all the prime ideals of A and it is also the intersection of the minimal prime ideals of A.

Clearly, if \(x \in A\) is nilpotent, it is contained in every prime ideal of A ( \(\S 1\) , no. 1, Definition 1). Conversely, let x be a non-nilpotent element of A; the set S of \(x^{k}\) (k an integer \(\geqslant 0\) ) is then a multiplicative subset of A not containing 0 and hence there exists a prime ideal p of A not meeting S (no. 5, Corollary 2 to Proposition 11) and a fortiori \(x \notin p\) ; this establishes the first assertion. To prove the second, it is sufficient to prove

LEMMA 2. Every prime ideal of a ring A contains a minimal prime ideal of A.

It is sufficient, by Zorn's Lemma, to show that the set P of prime ideals, ordered by the relation \(\supset\) , is inductive. Now, if G is a non-empty totally ordered subset of P, the intersection \(p_{0}\) of the ideals \(p \in G\) is also a prime ideal: for if \(x \notin p_{0}\) and \(y \notin p_{0}\) , there exists an ideal \(p \in G\) such that \(x \notin p\) and \(y \notin p\) , whence \(xy \notin p\) and a fortiori \(xy \notin p_{0}\) .

Remark. In § 4, no. 3, Corollary 3 to Proposition 14, we shall show that in a Noetherian ring the set of minimal prime ideals is finite; moreover we shall see later that every decreasing sequence of prime ideals in a Noetherian ring is stationary.

COROLLARY 1. The radical of an ideal a in a ring A is the intersection of the prime ideals containing a and it is also the intersection of the minimal elements of this set of prime ideals.

COROLLARY 2. For an ideal a of a ring A we denote by r(a) the radical of a. Then, for two ideals of a, b of A,

\[
\mathfrak {r} (a \cap b) = \mathfrak {r} (a b) = \mathfrak {r} (a) \cap \mathfrak {r} (b);
\]

in particular, if \(\mathfrak{a} \subset \mathfrak{b}\) , then \(\mathfrak{r}(\mathfrak{a}) \subset \mathfrak{r}(\mathfrak{b})\) .

For a prime ideal to contain \(\mathfrak{a} \cap \mathfrak{b}\) (or \(\mathfrak{ab}\) ), it is necessary and sufficient that it contains one of the ideals \(\mathfrak{a}, \mathfrak{b}\) ( \(\S 1\) , no. 1, Proposition 1).

PROPOSITION 14. For two ideals a, b of a ring to be relatively prime, it is necessary and sufficient that their radicals r(a) and r(b) be so.

The necessity of the condition is obvious since \(a \subset \mathfrak{r}(a)\) and \(b \subset \mathfrak{r}(b)\) ; the condition is sufficient by §1, no. 2, Proposition 3.

PROPOSITION 15. In a ring A let \(\mathfrak{a}\) be an ideal and \(\mathfrak{b}\) a finitely generated ideal contained in the radical of \(\mathfrak{a}\) . Then there exists an integer \(k > 0\) such that \(\mathfrak{b}^k \subset \mathfrak{a}\) .

Let \((b_{i})_{1\leqslant i\leqslant n}\) be a system of generators of b. By hypothesis there exists an integer h such that \(b_{i}^{h}\in\mathfrak{a}\) for \(1\leqslant i\leqslant n\) . If a product of nh elements, each of which is a linear combination of the \(b_{i}\) with coefficients in A, is expanded, each term is a multiple of a product of nh factors, each of which is equal to a \(b_{i}\) ; among these factors, at least h have the same index i for some i, hence, the product belongs to a and nh is the required integer k.

COROLLARY. In a Noetherian ring the nilradical is a nilpotent ideal.

PROPOSITION 16. Let B be a ring and A a subring of B. For every minimal prime ideal p of A there exists a minimal prime ideal q of B such that q ∩ A = p.

We write \(S = A - p\) ; then the ring \(A_p = S^{-1}A\) is identified with a subring of \(S^{-1}B\) (no. 1, Proposition 2) and on the other hand \(A_p\) has only a single prime ideal \(p'\) since \(p\) is minimal (no. 5, Proposition 11). As \(S^{-1}B\) is not reduced to 0 (since it contains \(A_p\) ), it has at least one prime ideal \(r'\) and therefore \(r' \cap A_p = p'\) ; if \(j = i_B^S\) and we write \(r = j^{-1}(r')\) , then

\[
i _ {\mathbf {A}} ^ {\mathbf {S}} (\mathfrak {r} \cap \mathbf {A}) \subset \mathfrak {r} ^ {\prime} \cap \mathbf {A} _ {p} = p ^ {\prime}
\]

hence \(\mathfrak{r} \cap \mathbf{A} \subset \mathfrak{p}\) and, as \(\mathfrak{p}\) is minimal, \(\mathfrak{r} \cap \mathbf{A} = \mathfrak{p}\) ; moreover, \(\mathfrak{r}\) is prime in \(\mathbf{B}\) ; if \(\mathfrak{q}\) is a minimal prime ideal of \(\mathbf{B}\) contained in \(\mathfrak{r}\) (Lemma 1), then a fortiori \(\mathfrak{q} \cap \mathbf{A} = \mathfrak{p}\) since \(\mathfrak{p}\) is minimal.

DEFINITION 5. A ring A is called reduced if its nilradical is reduced to 0, in other words if no element \(\neq 0\) of A is nilpotent.

If N is the nilradical of a ring A, A/N is reduced, for if the class mod. N of an element \(x \in A\) is nilpotent in A/N, this means that \(x^{h} \in N\) for some integer h, hence \(x^{hk} = 0\) for some integer k and \(x \in N\) .

PROPOSITION 17. Let A be a ring and \(\mathfrak{A}\) its nilradical. For every multiplicative subset S of A, \(S^{-1}\mathfrak{A}\) is the nilradical of \(S^{-1}A\) . In particular, if A is reduced, then \(S^{-1}A\) is reduced.

If \(x \in \mathbf{A}\) , \(s \in \mathbf{S}\) satisfy \((x / s)^n = x^n / s^n = 0\) , there exists \(s' \in \mathbf{S}\) such that \(s' x^n = 0\) (no. 1, Remark 3) and a fortiori \((s' x)^n = 0\) , hence \(s' x \in \mathfrak{N}\) and \(x / s = s' x / s' s \in \mathbf{S}^{-1} \mathfrak{N}\) ; the converse is immediate.

